% Standalone section fragment. Uses only macros/environments already in v0.1.
% Source claim identifiers are resolved in period_lattice_provenance.json.
\section{From finite periods to one integral target}\label{sec:period-lattice}

\status{The descent, gluing, integrality, and norm arguments below are proved here. Their inputs concerning the actual theta kernels and local divisibility are identified explicitly. For coefficient maps, ``rational'' means defined over the coefficient number field $K$.}

\subsection{The objects behind the maps}

The finite models are spaces of equivariant functions on decorated hermitian lattices. A point $d$ consists of a unimodular lattice in a positive definite hermitian space over a CM field, together with specified nonzero residue functionals and residue frames. A rational unitary group, together with a finite level subgroup, acts on these points. Once finiteness of the orbit set and stabilizers has been established, choose representatives $d_c$ and put $\Gamma_c=\operatorname{Stab}(d_c)$. For a coefficient representation $(V,\rho)$,
\[
 \mathcal A=\{f:f(gd)=\rho(g)f(d)\}
       \simeq\bigoplus_c V^{\Gamma_c},\qquad
 \mathcal L_{\Delta}=\{f\in\mathcal A:f(d)\in\Delta(d)\text{ for every }d\}.
\]
Here $\Delta(d)$ and its larger companion $\nabla(d)$ are explicit integral highest-weight modules; their polynomial definitions are recorded in the source notes. The initial source $C_u$ is the joint kernel of the residue-unit averaging operators
\[
 \operatorname{Nm}_P=\sum_{a\in\kappa(P)^\times}\langle a\rangle,
 \qquad C_u=\bigcap_{P\in S_{T,u}}\ker\operatorname{Nm}_P,
 \qquad M_u^{\rm ram}=C_u\cap\mathcal L_{\Delta}.
\]
Here $S_{T,u}$ is the set of primes dividing the chosen element $\pi_T(u)$, rather than necessarily every prime above the rational type prime $T_u$. The working lattice $N_u$ is a further arithmetic choice in $C_u$; it must not be identified with $M_u^{\rm ram}$. In the final divisibility bridge it is exactly the coordinate pullback of the record's $\texttt{NuLat}$. Its required local conditions remain inputs here. \textsf{[PL1, PL2, PL4, PL12]}

A kernel acts by the finite formula
\[
 (T_kf)(d')=\sum_c\frac{k(d_c,d')f(d_c)}{|\Gamma_c|}.
\]
Let $\mathcal T_{u,t}$ denote the chain of these operators, including the final evaluation or coupled-kernel step. The period in the code is
\begin{equation}\label{eq:pl-period}
 P_{u,t}(\phi,F)=\sum_c\frac{
 b_{\rm alg}\bigl(\operatorname{br}_u(\mathcal T_{u,t}\phi)(d_c),F(d_c)\bigr)}{|\Gamma_c|}.
\end{equation}
The coefficient pairing is bilinear: for polynomial coefficients it is a fixed normalization of $\sum_\alpha\alpha!f_\alpha g_\alpha$. Thus \eqref{eq:pl-period} is a $K$-bilinear finite sum, with visible stabilizer denominators. Constructing and estimating the individual theta kernels is additional work. \textsf{[PL3, PL15]}

On the row side, finite-group character and level projectors select a subspace $W_t$ of the analogous function space. Its lattice is
$M_t=W_t\cap\mathcal L_{\nabla}$. The selected theta-adjoint image gives an orthogonal quotient $\pi_t:W_t\twoheadrightarrow Q_t$, with its quotient metric. Its lattice is the \emph{image} $\pi_tM_t$. Rationality and common-image theorems for the theta maps produce isomorphisms
\[
 \iota_t:Q_t\overset\sim\longrightarrow Q:=Q_{t_0},
 \qquad \iota_{t_0}=1.
\]
Concretely, if $J_t$ is the injective complex theta map on $Q_t$ and all $J_t$ have the same image, then $\iota_t=J_{t_0}^{-1}J_t$; the assertion that this matrix lies in $K$ is a substantive input. \textsf{[PL5, PL6]}

\subsection{Canonical descent and comparison scalars}

\begin{proposition}[Descent of the period]\label{prop:pl-descent}
Suppose $P_{u,t_0}(\phi,F)=0$ whenever $\pi_{t_0}F=0$. There is a unique $K$-bilinear map $B_u:C_u\times Q\to K$ satisfying
\begin{equation}\label{eq:pl-base}
 B_u(\phi,\pi_{t_0}F)=P_{u,t_0}(\phi,F).
\end{equation}
It can be computed from any $K$-linear right inverse $r$ of $\pi_{t_0}$ by $B_u(\phi,y)=P_{u,t_0}(\phi,ry)$.

Fix an embedding $\sigma:K\hookrightarrow\mathbb C$. Suppose nonzero scalars $\xi_{u,t}\in\mathbb C$, with $\xi_{u,t_0}=1$, satisfy
\begin{equation}\label{eq:pl-complex}
 \xi_{u,t}\,\sigma(P_{u,t}(\phi,F))
     =\sigma\bigl(B_u(\phi,\iota_t\pi_tF)\bigr).
\end{equation}
Suppose also that whenever $\xi_{u,\mu(w)}=\sigma(a)$ for $a\in K$, one has $\operatorname{ord}_w(a)=0$. Then there are $\eta_{u,t}\in K^\times$, with $\eta_{u,t_0}=1$, such that
\begin{equation}\label{eq:pl-factor}
 P_{u,t}(\phi,F)=\eta_{u,t}B_u(\phi,\iota_t\pi_tF),
 \qquad \operatorname{ord}_w\eta_{u,\mu(w)}=0.
\end{equation}
\end{proposition}

\begin{proof}
Since $F-r\pi_{t_0}F$ is in the kernel of $\pi_{t_0}$, the displayed definition gives \eqref{eq:pl-base}. The same observation proves independence of $r$, and surjectivity of $\pi_{t_0}$ proves uniqueness.

If $\xi_{u,t}=\sigma(a)$, choose $\eta_{u,t}=a^{-1}$. Otherwise every value of $P_{u,t}$ must be zero: a nonzero value in \eqref{eq:pl-complex} would express $\xi_{u,t}$ as the image under $\sigma$ of the quotient of two elements of $K$. The right side is then zero too, so choose $\eta_{u,t}=1$. These choices give all the claims, including the finite-place unit condition. No nonvanishing assumption on $B_u$ was used.
\end{proof}

This is the source's actual right-inverse construction and scalar descent. The scalar is inverted when passing to \eqref{eq:pl-factor}. The kernel-vanishing input follows from the theta identity once $\ker\pi_t\subseteq\ker\Theta_t$ is known. The underlying linear algebra is short: if $\Theta_tW_t\subseteq E_0$, $\Theta_t^*E_0\subseteq W_t$, and $F\in W_t$ is orthogonal to $\Theta_t^*E_0$, then
$\|\Theta_tF\|^2=\langle F,\Theta_t^*\Theta_tF\rangle=0$.
Rationality and the theta identity still have to be supplied for the actual models. \textsf{[PL7, PL8, PL18]}

\subsection{Gluing the common lattice and using divisibility}

\begin{theorem}[One integral target for every word]\label{thm:pl-integral}
Let $M_t\subset W_t$ and $N_u\subset C_u$ be full $\OK$-lattices, and let $\pi_t$ and $\iota_t$ be as above. For each finite place $w$, choose a row parameter $\mu(w)$, equal to $t_0$ outside a finite set. Put
\begin{equation}\label{eq:pl-R}
 R=\bigcap_w\bigl(\iota_{\mu(w)}\pi_{\mu(w)}M_{\mu(w)}\bigr)_w
       \subset Q,
\end{equation}
where the subscript means localization over $\mathcal O_{K,w}$, inside the $K$-vector space. Then $R$ is a full lattice, uniquely determined by
$R_w=(\iota_{\mu(w)}\pi_{\mu(w)}M_{\mu(w)})_w$.

Assume \eqref{eq:pl-factor}. Let $\sigma_u(w)$ be finitely supported real numbers such that, for $\phi\in N_u$ and $F\in M_{\mu(w)}$,
\begin{equation}\label{eq:pl-val}
 P_{u,\mu(w)}(\phi,F)=0
 \quad\text{or}\quad
 \operatorname{ord}_wP_{u,\mu(w)}(\phi,F)\ge\sigma_u(w).
\end{equation}
Define
\begin{equation}\label{eq:pl-shift}
 k_u(w)=\lfloor\sigma_u(w)\rfloor,\qquad
 I_u=\prod_w\mathfrak p_w^{-k_u(w)},\qquad N'_u=I_uN_u.
\end{equation}
Then
\begin{equation}\label{eq:pl-integral}
 B_u(N'_u,R)\subseteq\OK,\qquad
 B_u:N'_u\longrightarrow R^\vee:=\operatorname{Hom}_{\OK}(R,\OK).
\end{equation}
The same $R$ works for every $u$.
\end{theorem}

\begin{proof}
Write $L_w=\iota_{\mu(w)}\pi_{\mu(w)}M_{\mu(w)}$ and $L_0=\pi_{t_0}M_{t_0}$. Images under surjective linear maps are full lattices, so all these lattices are full. They equal $L_0$ outside finitely many $w$.

Choose $0\ne d\in\OK$ with $dL_0\subseteq L_w\subseteq d^{-1}L_0$ for the exceptional $w$, and with $d$ in each exceptional prime. The finite torsion module $A=d^{-1}L_0/dL_0$ is the direct sum of its prime-primary components. In the component at each prime dividing $d$, prescribe the submodule given by $(L_w)_w/(dL_0)_w$. Pull back their direct sum to $d^{-1}L_0$. The resulting full lattice has exactly the required localizations; at primes not dividing $d$, both bounds localize to $L_0$.

A lattice is the intersection of its localizations: if $x$ fails to belong, the proper ideal $\{a\in\OK:ax\in L\}$ lies in a maximal ideal $\mathfrak p$, which prevents $x\in L_{\mathfrak p}$. This proves that the constructed lattice is \eqref{eq:pl-R}, as well as uniqueness.

For integrality, \eqref{eq:pl-factor} and the unit condition transfer \eqref{eq:pl-val} to $B_u$ on $N_u\times L_w$. The same bound holds after localization, since the new denominators are $w$-units. As valuations are integers, the bound is at least $\lceil\sigma_u(w)\rceil$. Multiplication by $I_u$ subtracts $k_u(w)$ from this bound. Thus $B_u(N'_u,R)$ is integral at every finite place, hence belongs to $\OK=\bigcap_w\mathcal O_{K,w}$. Finally, \eqref{eq:pl-R} contains no $u$.
\end{proof}

The proof also allows any finitely supported integer $k_u(w)\le\lceil\sigma_u(w)\rceil$; the released construction uses the floor. Its $\sigma_u(w)$ is the ramification index times the normalized gain recorded at $w$. A negative gain shrinks the lattice in \eqref{eq:pl-shift}, so ``shifted'' is more accurate than ``enlarged.'' \textsf{[PL9--PL12]}

\subsection{Why the quotient also gives the right norm bound}

At an infinite place $v$, write the base period as
$P_{u,t_0}(\phi,F)=P_{\rm end}(T_r\cdots T_1\phi,F)$.
Give $Q_{t_0}$ the quotient metric from $W_{t_0}$. Suppose
\[
 |P_{\rm end}(z,F)|\le p_v\|z\|\|F\|,\quad
 \|T_j\|\le a_{j,v},\quad
 \|\iota_{t_0}q\|_Q^2\ge s_v\|q\|_{Q_{t_0}}^2,
 \qquad p_v,a_{j,v},s_v>0.
\]
Here the metric on the common target $Q$ can differ from the quotient metric, even though $\iota_{t_0}$ is the identity in the chosen coordinates. For $y\in Q_v$, take the orthogonal, minimum-norm lift $F$ of $\iota_{t_0}^{-1}y$. Substituting this $F$ in \eqref{eq:pl-base} proves
\begin{equation}\label{eq:pl-norm}
 \|B_u\|_v^2\le\frac{p_v^2\prod_j a_{j,v}^2}{s_v}.
\end{equation}
There is no loss from descent through the quotient itself: a functional vanishing on its kernel has exactly the same norm as the descended functional.

In the source, $p$ ranges over a finite nonempty set of spectral pieces,
$\kappa_v>0$, and every $\lambda_{j,p,v}>0$. The eigenvalue comparison gives
$s_v=\kappa_v\min_p\prod_j\lambda_{j,p,v}>0$.
Consequently, if $B_u\ne0$, its squared-norm height satisfies
\begin{equation}\label{eq:pl-height}
 h(B_u)\le\operatorname{mean}_v\!\left(
 \log p_v^2+\sum_j\log a_{j,v}^2-\log\kappa_v
       -\min_p\sum_j\log\lambda_{j,p,v}\right).
\end{equation}
The mean is over all infinite places, with the source's equal weights for totally complex $K$. Formula \eqref{eq:pl-norm} also applies when $B_u=0$; then one uses positive majorants without taking $\log0$. \textsf{[PL11, PL20]}

What remains is now localized: prove the rational common-image and theta comparison identities; establish \eqref{eq:pl-val} for the chosen $N_u$; and bound the chain norms and eigenvalues at every infinite place. This section supplies the complete algebraic passage from those inputs to integral maps with a common target and the required form of height estimate. It does not yet establish their quantitative arithmetic inputs.
